
\medskip
\phantomsection\label{prop:sga3-VIB-6.5.1}%
\noindent\textbf{Proposition 6.5.1.---}
\emph{Let \(G\) be an \(S\)-group locally of finite presentation
and quasi-separated over \(S\), let \(H\) be a smooth \(S\)-group
with connected fibers, and let \(u:H\to G\) be a monomorphism.
Let \(N\) be the normalizer of \(H\) in \(G\)
(cf.\ I or \hyperref[def:sga3-VIB-6.1]{6.1}). We shall see
\textup{(}XI, \SGARefLink{sga3:xi:corollary:6:11}{6.11}\textup{)}%
\begingroup
\renewcommand{\thefootnote}{\ensuremath{(\ast)}}%
\footnote[1]{\label{note:sga3-VIB-6.5.1-star}See the note \((\ast)\) to the proof of
\hyperref[note:sga3-VIB-6.2-star]{6.2}; the 1970 text refers to it as the note on page~40, the page of this Expos\'e
on which it stands.}%
\endgroup{}
that \(N\) is representable
by a closed group subscheme of \(G\), of finite presentation over
\(G\). Under these circumstances, the following conditions are
equivalent:}
\begin{enumerate}
\item[\textup{\SGAInlineRefTarget{prop:sga3-VIB-6.5.1:item:i}(i)}]\phantomsection\label{cond:sga3-VIB-6.5.1-i}
  \emph{The canonical morphism \(H\to N\) is an open immersion.}
\item[\textup{\SGAInlineRefTarget{prop:sga3-VIB-6.5.1:item:ii}(ii)}]\phantomsection\label{cond:sga3-VIB-6.5.1-ii}
  \emph{\(N^0=H\)
  (cf.\ \hyperref[thm:sga3-VIB-3.10]{3.10}).}
\item[\textup{\SGAInlineRefTarget{prop:sga3-VIB-6.5.1:item:iii}(iii)}]\phantomsection\label{cond:sga3-VIB-6.5.1-iii}
  \emph{For every \(s\in S\), one has
  \(H_s=(N_s)^0\).}
\end{enumerate}

Condition \textup{\SGARefLink{prop:sga3-VIB-6.5.1:item:i}{(i)}} implies \textup{\SGARefLink{prop:sga3-VIB-6.5.1:item:ii}{(ii)}} by
\hyperref[lem:sga3-VIB-3.10.1]{Lemma 3.10.1}, since \(H^0=H\).
It is also clear that \textup{\SGARefLink{prop:sga3-VIB-6.5.1:item:ii}{(ii)}} implies \textup{\SGARefLink{prop:sga3-VIB-6.5.1:item:iii}{(iii)}}, for
\[
  H_s=(N^0)_s=(N_s)^0.
\]

Finally, let us prove that \textup{\SGARefLink{prop:sga3-VIB-6.5.1:item:iii}{(iii)}} implies \textup{\SGARefLink{prop:sga3-VIB-6.5.1:item:i}{(i)}}.
Since \(H_s=(N^0)_s\) for every \(s\in S\), the morphism
\[
  H_s\longrightarrow N_s
\]
is an open immersion. Moreover, \(H\) and \(N\) are locally of
finite presentation over \(S\), and \(H\) is flat over \(S\).
Consequently, \(H\to N\) is an open immersion
(EGA~IV\(_4\), 17.9.5).\footnote{Translator's note: the Gille--Polo re-edition (2011), Expos\'e~VI\textsubscript{B},
inserts after 6.5.1, as its 6.5.2--6.5.5, Expos\'e~XI, \SGARefLink{sga3:xi:theorem:6:8}{6.8}--6.11 (in this
edition in Expos\'e~XI), as A.~Grothendieck suggests at the beginning of XI, \S\,6, and refers to its 6.5.5
instead of XI, 6.11. In XI, 6.8 it assumes the fibres of \(X\) geometrically irreducible instead of geometrically
connected; for \(X\) smooth over \(S\) the two conditions agree, a connected normal noetherian scheme being
irreducible. It notes that, by M.~Raynaud and L.~Gruson, \emph{Crit\`eres de platitude et de projectivit\'e},
Invent.\ Math.\ \textbf{13} (1971), I, 3.3.4~(iii) and 4.1.1, it is enough there that \(X\) be flat over \(S\),
with geometrically irreducible fibres without embedded components; it also expands the proofs, using
EGA~IV\(_3\), 11.10.10 for the schematic density of an open subscheme containing the image of the section.}

